\section{Results}

\begin{table*}[t]
\centering
\caption{The 13 questions, their category and hub scope, and each system's
score. Baselines: one answer per slice (A, B, C) for each of the three
traces, so 9 per question. \sysname: one answer per trace, so 3 per
question. Correct = 1, partly correct = \textonehalf, wrong = 0.}
\label{tab:questions}
\small
\begin{tabular}{@{}rp{3.05in}llrrr@{}}
\toprule
Q & Question & Category & Hub scope & LLM (/9) & RAG (/9) & \sysname (/3) \\
\midrule
1  & What does this trace record? & Comprehension & Multi & 3 & 2.5 & 3 \\
2  & What does each thread of \texttt{batch\_server} do? & Comprehension & Multi & 3 & 7.5 & 3 \\
3  & Which CPU(s) does \texttt{batch\_server} run on? & Comprehension & Single & 6 & 7.5 & 3 \\
4  & Which thread is the bottleneck of \texttt{batch\_server}? & Root cause & Multi & 2 & 7.5 & 3 \\
5  & Why is \texttt{worker-2} slow? & Dependency & Multi & 1 & 7.5 & 3 \\
6  & Is \texttt{worker-2} slow because it uses more CPU than the other workers? & Dependency & Multi & 4 & 5.5 & 3 \\
7  & What is \texttt{worker-2} waiting on, and which thread wakes it up? & Dependency & Single & 5.5 & 6.5 & 3 \\
8  & Why does \texttt{flusher} hold the lock for so long? & Dependency & Single & 1 & 7.5 & 3 \\
9  & Break down how \texttt{worker-2} spends its time in a typical round. & Quantitative & Single & 2 & 2 & 3 \\
10 & What change would fix the slowdown, and how much faster would rounds be? & Remediation & Multi & 0 & 7.5 & 3 \\
11 & In a typical round, how much later does \texttt{worker-2} finish than the other workers? & Quantitative & Multi & 2 & 5.5 & 3 \\
12 & How could \texttt{flusher} hold \texttt{cache\_lock} for less time, and how long afterwards? & Remediation & Single & 6.5 & 7.5 & 3 \\
13 & Does another thread hold up \texttt{flusher}, or does it delay itself? & Root cause & Single & 5.5 & 5.5 & 3 \\
\midrule
   & \textbf{Total} & & & 41.5/117 & 80/117 & 39/39 \\
   & & & & (35.5\%) & (68.4\%) & (100\%) \\
\bottomrule
\end{tabular}
\end{table*}

\begin{table}[t]
\centering
\caption{Scores as a percentage of the maximum, by slice, category and hub
scope. \sysname reads the whole trace, so it has no slices.}
\label{tab:breakdown}
\small
\begin{tabular}{@{}lrrr@{}}
\toprule
 & LLM & RAG & \sysname \\
\midrule
Slice A (head)      & 4  & 29 & -- \\
Slice B (middle)    & 55 & 90 & -- \\
Slice C (sampled)   & 47 & 86 & -- \\
\midrule
Comprehension (Q1--3)          & 44 & 65 & 100 \\
Root cause (Q4, 13)            & 42 & 72 & 100 \\
Dependency (Q5--8)             & 32 & 75 & 100 \\
Quantitative (Q9, 11)          & 22 & 42 & 100 \\
Remediation (Q10, 12)          & 36 & 83 & 100 \\
\midrule
Single-hub (6 questions)       & 49 & 68 & 100 \\
Multi-hub (7 questions)        & 24 & 69 & 100 \\
\midrule
All                            & 35.5 & 68.4 & 100 \\
\bottomrule
\end{tabular}
\end{table}

\textbf{Questions and grading.} The 13 questions
(Table~\ref{tab:questions}) fall into five categories: comprehension (what
the trace and the threads do), root cause, dependency (what waits on what),
quantitative analysis and remediation. Following TAAF~\cite{ezaz2026taaf},
we also label each question by \emph{hub scope}: single-hub if one central
entity and its direct relationships are enough to answer it, such as one
thread's waits, and multi-hub if the answer combines or compares evidence
across several entities, such as the wait chain through three threads.
TAAF found multi-hub questions harder for LLMs; the set has 6 single-hub
and 7 multi-hub questions. Each answer is graded against the expected
answer and the measured values from the program's round log.

\textbf{Overall.} \sysname answered all 13 questions correctly on all
three traces, against 35.5\% for the LLM alone and 68.4\% with RAG. Trace
length made little difference to any system: the baselines scored
32--41\% and 64--71\% per trace, and \sysname 13 of 13 on each.

\textbf{The baselines depend on where they look.} Slice A, what a user
gets by pasting a trace into a chat, ends before the program starts, so
the LLM alone scores 4\% on it; RAG's 29\% comes from repeating the
program description. With slices B and C, which contain the program's
activity, RAG reaches 86--90\% (Table~\ref{tab:breakdown}): the
description tells the model what to look for, and the slice confirms it. But
choosing such a slice requires already knowing where in the trace the
program runs. \sysname reads the whole trace and needs no such choice.

\textbf{Multi-hub questions.} The LLM alone falls from 49\% on single-hub
questions to 24\% on multi-hub ones, consistent with TAAF. RAG closes the
gap (68\% and 69\%) because the program description states how the
threads relate. \sysname scores 100\% on both, because its process
explainer measures the wait chain across threads rather than taking it
from a document.

\textbf{Quantitative answers.} Quantitative questions are the weakest
category for both baselines (22\% and 42\%). A slice holds at most a few
rounds, so a baseline measures one round: RAG's answers put
\texttt{worker-2}'s lock wait at 24.86 or 27.06\,ms, against a mean of
24.7\,ms over all rounds, and only two of nine time breakdowns (Q9) were
complete. \sysname's figures are averages over every round and match the
ground truth closely: the per-round gap between \texttt{worker-2} and the
other workers (Q11) is 22.173, 22.127 and 22.156\,ms on the three traces,
against 22.173, 22.128 and 22.156\,ms measured by the program, and its
lock wait is within 0.010\,ms of the measured mean.

\textbf{Retrieval.} With the question as the only query, RAG did not
retrieve the program description for Q1 or Q13, so it scored no better
than the LLM alone on them. \sysname's queries draw on its findings, and
it retrieved the description for Q1.

\textbf{Limitations.} The evaluation covers one program with one bug, and
\sysname ran once per trace. Its thread questions carried thread IDs,
which the baselines' did not, and its answers within a batch could see
earlier ones. One \sysname answer (Q6, 30-minute trace) also reported two
``abnormally late'' cycles that the ground truth does not show; the claim
came from the auxiliary statistical module and did not change the answer.
All answers were graded by one author.

\textbf{Summary.} Given only trace text, an LLM answered about a third of
the questions correctly, and adding the program's documentation raised
this to two thirds, but only when its slice happened to contain the
program's activity, and its numbers came from single rounds. Grounded in
findings computed from the whole trace, \sysname answered all 13 questions
correctly on all three traces, with its key measurements within 0.01\,ms
of the ground truth.
